\section{Introduction}
\label{sec:introduction}

\subsection{Background and motivation}

Thermal energy storage increases the flexibility of district heating systems by decoupling heat production from demand~\cite{guelpa2019review,vandermeulen2018controlling}. In systems with combined heat and power (CHP) plants, storage allows heat production to follow electricity prices, and pit thermal energy storage (PTES) extends this flexibility from hours to days or weeks at low specific cost~\cite{dahash2019seasonal,schmidt2018design}. Energy-system studies attribute a substantial value to such storage. For Denmark and its neighbouring countries, \citet{sifnaios2023impact} find that large-scale thermal storage allows higher wind and solar capacities and less curtailment, and \citet{christensen2024role} find that it reduces fuel consumption and curtailment during the transition to a fully decarbonised system. For Europe, \citet{burkhardt2026flexibility} show that storage enables a deeper electrification of district heating and that PTES is operated in cycles of about a month rather than seasonally. For Germany, \citet{schauss2026mind} estimate that PTES lowers annual system costs by 135--345~million~EUR compared with tank storage alone, even when the temperature difference between the storage and the network is taken into account. All of these studies determine the operation of the storage from the perspective of the whole system.

Whether this value is realised in operation depends on who operates the storage and on the prices that this party faces. The Greater Copenhagen district heating system illustrates the point. Production from CHP plants, waste incineration, heat pumps and peak boilers is dispatched centrally by Varmelast to minimise system cost, transmission companies such as VEKS buy and transport the heat, and local distribution companies supply end users. The dispatch determines quantities, whereas heat is settled at prices agreed in contracts, so the hourly value of heat that the dispatch implies is not passed on~\cite{frolke2022market}. The Høje Taastrup PTES (3.3~GWh) is jointly owned by VEKS and the local distribution company Høje Taastrup Fjernvarme (HTF); it is charged from the transmission network, scheduled by Varmelast and discharged into HTF's network. In 2024, its measured discharge corresponds to about nine equivalent full cycles (\cref{sec:results}), compared with a design range of 25--30 cycles per year~\cite{sifnaios2025performance}. Several factors contribute, including network restrictions and low seasonal demand~\cite{sifnaios2025performance}. One factor that can be addressed through contract design is the settlement: HTF pays the same monthly tariff for direct and discharged heat, and discharge adds pumping cost, so implementing a discharge schedule offers no financial benefit to HTF even when it reduces system cost.

\subsection{Related work}

\paragraph{Marginal-cost and dynamic heat pricing}
Heat prices based on marginal cost have long been proposed as an alternative to cost-plus tariffs that are fixed for a year or a season~\cite{sjoedin2004marginal,difs2009pricing,zhang2013equivalent,sun2016marginal}; \citet{li2015review} review the pricing mechanisms in use. \citet{dominkovic2018influence} simulate district heating supply and demand in the manner of a day-ahead electricity market, in which the hourly heat price is set where the marginal-cost supply curve meets demand, and show for Sønderborg and Espoo how waste heat and thermal storage change this price and that dynamic pricing favours the feed-in of waste heat. Wholesale heat markets with hourly prices have been modelled for the Netherlands~\cite{liu2019marginal} and for the integration of excess heat~\cite{doracic2021utilizing}. In these studies the price is the marginal value of heat, and it is paid to producers or charged to consumers.

\paragraph{Price signals for flexibility}
A second group of studies uses time-varying heat prices to activate flexibility on the demand side. \citet{cai2018demand} use the marginal cost of heat in Copenhagen as a price signal for demand-side management in an urban network. \citet{bai2022variable} formulate variable heat pricing as a leader--follower problem in which the heat provider announces hourly prices day-ahead and consumers adjust their demand. \citet{mokhtari2025dynamic} argue that a price based on marginal cost alone does not steer demand response adequately and design day-ahead price profiles by inverse optimisation on a characterisation of the buildings' flexibility. This characterisation, the Flexibility Function, describes demand as a dynamic response to a price signal~\cite{junker2018characterizing,junker2020stochastic}. The Frigg framework~\cite{schledorn2022a,schledorn2024frigg} soft-links such a response model with an energy-system dispatch model: instead of setting demand directly, the central problem chooses a bounded price signal, and the response to this signal is evaluated by the separate model, after which the dispatch is solved again with the response fixed. \citet{dominkovic2020implementing} soft-link a planning model with a short-term operational model through a characterisation of short-term flexibility. In an application of Frigg to Denmark, heat storage yields larger savings than demand response alone~\cite{schledorn2024frigg}, but the storage is operated by the central model. With regard to practice, \citet{selvakkumaran2021business} find that the price models of district energy companies rarely express the value of flexibility to customers, and \citet{lygnerud2024capturing} test price models for integrating heat pumps into district heating networks; demand response in district heating and the barriers to flexibility at the interface with the electricity system are reviewed in~\cite{guelpa2021demand,sneum2021barriers}.

\paragraph{Coordination between actors}
District heating systems with several companies raise the question of how the actors are coordinated. Third-party access to networks has been analysed with regard to competition and to the integration of independent producers~\cite{soederholm2011market,buerger2019third}. For Copenhagen, \citet{frolke2022market} compare two ways of integrating excess-heat producers: participation in the heat market, and self-scheduling in response to a published price signal. They find that a simple signal is sufficient while the responding capacity is small and becomes increasingly inefficient as it grows. Market mechanisms that account for the network~\cite{frolke2022network} and for the sequential clearing of heat and electricity markets~\cite{mitridati2020heat} have been proposed. In electricity systems, the corresponding question is how flexibility connected at distribution level can be used for the needs of the transmission system, for which several coordination schemes between the system operators have been defined~\cite{gerard2018coordination}.

\subsection{Research gap and contributions}

Three observations follow from this literature. First, studies of heat pricing address the price of heat for consumers or producers and use the marginal cost of heat as the signal. Second, frameworks for price-based flexibility address demand response, while storage is operated centrally. Third, energy-system studies of PTES assume an operation that is optimal for the system. The settlement of a storage that is scheduled by one company and operated by another, at the boundary between a transmission and a distribution company, has to our knowledge not been addressed. In this setting the marginal cost of heat is not the relevant signal: discharging replaces direct supply, so what matters is the difference between the value of delivering heat now and the value of keeping it stored, and the signal has to hold against the operating cost and the local production of the company that responds to it.

This paper addresses the question: \emph{How can a dynamic PTES discharge tariff be derived from the value of storage in coordinated heat scheduling, and can it make both the local operator and the upstream side better off than the current settlement?} The approach follows the principle of indirect control in Frigg~\cite{schledorn2022a,schledorn2024frigg}: the response is not prescribed but results from a price, and its effect on the system is evaluated by solving the dispatch again around it. It differs in three respects. The price is not optimised as a variable but derived from the dual values of the coordinated schedule, in which the hours are linked by the storage, so that the heat balance and the storage each have a marginal value and the credit is their difference. The responding party is a company that minimises its own expenditure with a storage, pumps and heat pumps. And because the price is an actual payment, the system cost and the accounts of both parties are reported separately.

The contributions are (i) a derivation of an hourly discharge credit from the marginal values of a coordinated heat-system schedule, (ii) an evaluation of its effect on system operating cost and on each stakeholder for 2024, relative to the measured operation of the storage, (iii) an analysis of how the storage value can be shared between the local operator and the upstream side, and of which inefficiencies a discharge tariff alone cannot address, and (iv) an assessment of the value of implementing the schedule exactly and of using the design discharge capacity of 30~MW instead of the current limit of about 23~MW.

\Cref{sec:case_study} describes the case study, \cref{sec:methodology} the model and tariff derivation, \cref{sec:results} the results, \cref{sec:discussion} their implications and limitations, and \cref{sec:conclusion} concludes.
